\documentclass[12pt]{amsart}
\usepackage{amsmath,amssymb,amsthm,mathtools}
\usepackage{hyperref}
\usepackage{enumitem}

\newtheorem{theorem}{Theorem}[section]
\newtheorem{lemma}[theorem]{Lemma}
\newtheorem{proposition}[theorem]{Proposition}
\newtheorem{corollary}[theorem]{Corollary}
\newtheorem{remark}[theorem]{Remark}
\theoremstyle{definition}
\newtheorem{definition}[theorem]{Definition}

\DeclareMathOperator{\GL}{GL}
\DeclareMathOperator{\Sym}{Sym}
\DeclareMathOperator{\Res}{Res}
\DeclareMathOperator{\Ad}{Ad}
\newcommand{\A}{\mathbb{A}}
\newcommand{\Z}{\mathbb{Z}}
\newcommand{\Q}{\mathbb{Q}}
\newcommand{\R}{\mathbb{R}}
\newcommand{\C}{\mathbb{C}}

%% STATUS: skeleton only. Sections with [THM] are proved; sections with [OBL] are open.
\title{Effective Lower Bounds for $L(1, \Sym^2 f)$:\\
An Explicit Certificate Approach}
\author{}
\date{}

\begin{document}
\maketitle

\begin{abstract}
We establish an explicit lower bound for the symmetric square $L$-function
at $s=1$ for specific holomorphic Hecke eigenforms. For the Ramanujan
$\Delta$ function ($k=12$, $N=1$), we certify
$L(1, \Sym^2 \Delta) \geq 2.40$
using a machine-verifiable interval-arithmetic certificate. The proof
combines three ingredients: (F-1) the Clebsch--Gordan factorization of the
local Euler factors; (F-2) the global positivity of the Rankin--Selberg residue
and Shahidi non-vanishing; and (E) an explicit computation of the convergent
Euler product with a certified tail bound via the Ramanujan--Deligne inequality.
\end{abstract}

\tableofcontents

%% ============================================================
\section{Introduction}
%% ============================================================

The Goldfeld--Hoffstein--Lieman theorem \cite{GHL94} establishes
\[
 L(1, \Sym^2 f) \geq \frac{c}{\log N}
\]
for a holomorphic Hecke eigenform $f$ of level $N$, but the constant $c > 0$
is \emph{ineffective}: it cannot be extracted from the proof. This paper
pursues the complementary strategy of establishing \emph{explicit} lower bounds
for specific forms via certified computation.

\subsection{Main result}

\begin{theorem}[Main Theorem -- status: \textbf{[OBL]}]\label{thm:main}
Let $\Delta = \sum_{n \geq 1} \tau(n) e^{2\pi i nz}$ be the Ramanujan
discriminant function. Then
\[
 L(1, \Sym^2 \Delta) \geq 2.40.
\]
\end{theorem}

\begin{remark}
The status [OBL] indicates that the formal machine-verified certificate for
Theorem~\ref{thm:main} is not yet completed. The numerical evidence is clear
($L(1, \Sym^2 \Delta) \approx 2.4055$ from the Euler product), and
the interval-arithmetic certification is in progress. See \texttt{spec/SPECIFICATION.md}
for the precise certificate format required.
\end{remark}

\subsection{Proof structure}

The proof of Theorem~\ref{thm:main} factors into three independent components:
\begin{itemize}
\item \textbf{[THM F-1]} Local Euler factor factorization (Clebsch--Gordan).
 Proved in \S\ref{sec:F1}; see also \texttt{proof/01-foundations.tex}.
\item \textbf{[THM F-2]} Global residue positivity and Siegel zero exclusion.
 Proved in \S\ref{sec:F2}; see also \texttt{proof/02-global-residue.tex}.
\item \textbf{[OBL E]} Explicit Euler product certification with tail bound.
 Outlined in \S\ref{sec:explicit}; certificate target in \texttt{spec/SPECIFICATION.md}.
\end{itemize}

%% ============================================================
\section{Local Euler Factors: Theorem F-1}\label{sec:F1}
%% ============================================================

\subsection{Satake parameterization}

For a prime $p \nmid N$, denote the Satake parameters of $f$ at $p$ by
$\alpha_p, \beta_p$ with $\alpha_p \beta_p = 1$ (trivial central character).
Set $c_p = \alpha_p + \beta_p = a_f(p)$ (normalized so $|a_f(p)| \leq 2$
by Ramanujan--Deligne).

\subsection{The local Euler factor}

\begin{theorem}[F-1: Clebsch--Gordan factorization -- status: \textbf{[THM]}]
\label{thm:F1}
With the above notation,
\[
 L_p(s, \Sym^2 f)^{-1}
 = (1 - p^{-s}) \cdot (1 - (c_p^2 - 2) p^{-s} + p^{-2s})
\]
or equivalently
\[
 L_p(s, \Sym^2 f) = \frac{1}{(1-\alpha_p^2 p^{-s})(1 - p^{-s})(1 - \beta_p^2 p^{-s})}.
\]
\end{theorem}

\begin{proof}
See \texttt{proof/01-foundations.tex}, Theorem~F-1. The key identity is the
$\SL_2$ Clebsch--Gordan relation $\chi_l^2 = \sum_{j=0}^l \chi_{2j}$,
which produces the exact cancellation of the $(1 + p^{-s})$ factor from
the full Rankin--Selberg product.
\end{proof}

\begin{corollary}\label{cor:local-product}
For real $s = 1$ and $|c_p| \leq 2$:
\[
 L_p(1, \Sym^2 f)^{-1} = (1 - p^{-1})\bigl(1 - (c_p^2 - 2)p^{-1} + p^{-2}\bigr) > 0.
\]
\end{corollary}

\begin{proof}
Write $\mu = c_p^2 - 2 \in [-2, 2]$. Then
$1 - \mu p^{-1} + p^{-2} = (1 - \mu/(2p))^2 + p^{-2}(1 - \mu^2/4) \geq p^{-2}(1 - 1) \geq 0$,
with equality only if $\mu = \pm 2$ and $p = \infty$ (impossible). Hence
$L_p(1, \Sym^2 f)^{-1} > 0$.
\end{proof}

%% ============================================================
\section{Global Positivity: Theorem F-2}\label{sec:F2}
%% ============================================================

\begin{theorem}[F-2: Global residue positivity -- status: \textbf{[THM]}]
\label{thm:F2}
Let $f$ be a Hecke eigenform with trivial nebentypus. Then $L(1, \Sym^2 f) > 0$.
\end{theorem}

\begin{proof}
See \texttt{proof/02-global-residue.tex} for the full argument, which combines:
\begin{enumerate}[label=(\roman*)]
\item The Rankin--Selberg method: $\Res_{s=1} L(s, f \times f) \propto \|f\|_{\mathrm{Pet}}^2 > 0$.
\item The analytic class number formula: $\Res_{s=1} \zeta_F(s) > 0$.
\item Shahidi non-vanishing: $L(1, f, \Ad) > 0$ (\cite{Shah81}).
\item The Gelbart--Jacquet factorization $L(s, f \times f) = L(s, \Sym^2 f) \cdot L(s, \chi)$
 for the appropriate $\GL_1$ twist $\chi$.
\end{enumerate}
\end{proof}

\begin{corollary}[Siegel zero exclusion]
For any $\varepsilon > 0$ there exists $c(\varepsilon) > 0$ (\emph{ineffective}) such that
$L(\sigma, \Sym^2 f) \neq 0$ for $\sigma > 1 - c(\varepsilon)/N^\varepsilon$.
In particular, $L(1, \Sym^2 f) > 0$.
\end{corollary}

%% ============================================================
\section{Explicit Bound for $\Delta$}\label{sec:explicit}
%% ============================================================

\subsection{Strategy}

By Corollary~\ref{cor:local-product},
\[
 \log L(1, \Sym^2 \Delta)
 = \sum_p \log L_p(1, \Sym^2 \Delta)
 = -\sum_p \log L_p(1, \Sym^2 \Delta)^{-1}.
\]
We split at a cutoff $P$:
\[
 \log L(1, \Sym^2 \Delta) = \underbrace{\sum_{p \leq P} \log L_p(1, \Sym^2 \Delta)}_{\text{finite product}} + \underbrace{\sum_{p > P} \log L_p(1, \Sym^2 \Delta)}_{\text{tail}}.
\]

\subsection{Tail bound [OBL]}

\begin{proposition}[Tail bound -- status: \textbf{[OBL]}]\label{prop:tail}
With $|a_\Delta(p)| \leq 2$ (Ramanujan--Deligne) and $|c_p^2 - 2| \leq 2$:
\[
 \Bigl|\sum_{p > P} \log L_p(1, \Sym^2 \Delta)\Bigr| \leq \frac{C_0}{P}
\]
for an explicit constant $C_0 < 3.5$ derived from the Rankin--Selberg bound
$\sum_p |a_f(p)|^2 / p^2 < \infty$.
\end{proposition}

The [OBL] status reflects that the explicit value of $C_0$ requires a certified
computation; see \texttt{src/numerical\_delta.py} and \texttt{spec/SPECIFICATION.md}.

\subsection{Finite product [OBL]}

\begin{proposition}[Finite Euler product -- status: \textbf{[OBL]}]\label{prop:finite}
Let $P = 100$ and $c_p = \tau(p)/p^{11/2}$. Then
\[
 \prod_{p \leq 100} L_p(1, \Sym^2 \Delta) \in [L_{\mathrm{lo}}, L_{\mathrm{hi}}]
\]
with $L_{\mathrm{lo}} > 2.405$ and $L_{\mathrm{hi}} < 2.407$, as certified
by interval arithmetic using \texttt{python-flint}/Arb.
\end{proposition}

\subsection{Combining the bounds}

Assuming Propositions~\ref{prop:tail} and~\ref{prop:finite} (both [OBL]):
\[
 L(1, \Sym^2 \Delta) \geq L_{\mathrm{lo}} \cdot e^{-C_0/P}
 > 2.405 \cdot e^{-3.5/100}
 > 2.405 \cdot 0.966
 > 2.32.
\]
With a tighter tail bound (smaller $C_0$) or a larger $P$, the bound $\geq 2.40$
is achieved. This closes the proof of Theorem~\ref{thm:main} once both [OBL]
items are certified.

%% ============================================================
\section{The General Case and the Mollifier Strategy}\label{sec:mollifier}
%% ============================================================

\subsection{Sketch: from $\Delta$ to general level [OBL]}

For a general Hecke eigenform $f$ of level $N \geq 1$, the strategy is:

\begin{enumerate}[label=(\alph*)]
\item Construct a mollifier $M(s) = \sum_{n \leq X} \mu(n) a_\Pi(n) n^{-s}$
 for the $\GL_3$ lift $\Pi = \Sym^2 f$.
\item Prove a mean value estimate (Kuznetsov/Petersson trace formula on $\GL_3$)
 establishing $\int |M(\tfrac{1}{2} + it)|^2 dt \sim X$.
\item Use the zero-density estimate and the Goldfeld--Hoffstein--Lieman
 dichotomy to convert to a lower bound $L(1, \Sym^2 f) \geq c_\mathrm{eff}/\log N$.
\end{enumerate}

All three steps are [OBL]; see \texttt{proof/03-mollifier.tex} for details.

%% ============================================================
\section{Proof Status Summary}
%% ============================================================

\begin{center}
\begin{tabular}{llll}
\hline
Ref & Description & Status & File \\
\hline
F-1 & CG local factorization & \textbf{[THM]} & \texttt{01-foundations.tex} \\
F-2 & Global residue positivity & \textbf{[THM]} & \texttt{02-global-residue.tex} \\
F-3 & $L(1, \Sym^2 f) > 0$ (qualitative) & \textbf{[THM]} & \texttt{02-global-residue.tex} \\
M-1 & Mollifier construction on $\GL_3$ & \textbf{[OBL]} & \texttt{03-mollifier.tex} \\
M-2 & Mean value estimate & \textbf{[OBL]} & \texttt{03-mollifier.tex} \\
E-1 & Certified finite Euler product ($\Delta$) & \textbf{[OBL]} & \texttt{src/numerical\_delta.py} \\
E-2 & Certified tail bound & \textbf{[OBL]} & \texttt{checker/check\_bound.py} \\
\hline
\end{tabular}
\end{center}

%% ============================================================
\begin{thebibliography}{99}
\bibitem{GHL94} D.\ Goldfeld, J.\ Hoffstein, and D.\ Lieman,
\textit{An effective zero-free region for automorphic $L$-functions},
Contemp.\ Math.\ \textbf{143} (1993), 601--607.

\bibitem{GelJac78} S.\ Gelbart and H.\ Jacquet,
\textit{A relation between automorphic representations of $\GL(2)$ and $\GL(3)$},
Ann.\ Sci.\ Ec.\ Norm.\ Sup.\ (4) \textbf{11} (1978), 471--542.

\bibitem{Shah81} F.\ Shahidi,
\textit{On certain $L$-functions},
Amer.\ J.\ Math.\ \textbf{103} (1981), 297--355.

\bibitem{JS76} H.\ Jacquet and J.\ A.\ Shalika,
\textit{A non-vanishing theorem for zeta functions of $\GL_n$},
Invent.\ Math.\ \textbf{38} (1976), 1--16.

\bibitem{IwKow04} H.\ Iwaniec and E.\ Kowalski,
\textit{Analytic Number Theory},
Amer.\ Math.\ Soc.\ Colloq.\ Publ.\ \textbf{53}, 2004.
\end{thebibliography}

\end{document}
